
\chapter{The Archimedean Principle}

If \(a\) and \(b\) are real numbers with \(a > 0\), then there exists \(n \in \N\) 
such that \(na > b\).

Equivalently, for any real number \( x \in \R \), there exists a natural number \( n \in \N \) such 
that \( n > x = \frac{b}{a} \). 

\textbf{Proof:}

We prove the Archimedean Principle by contradiction.

Assume that there exists some real number \( x \in \R \) such that \( n \leq x \) for all 
\( n \in \N \). That is, \(x\) is an upper bound for the set \( \N \subset \R\), baked up 
by the Completeness Axiom of the real numbers.

Let \( s = \sup(\N) \) be the least upper bound of \( \N \). Then, 

\[
	s - 1 < s \implies \exists m \in \N \text{ s. t. } m > s - 1
\]

But, if we add 1 to both sides we get  

\[
	s < m + 1 \in \N
\]

This is a contradiction to the assumption that \(s\) is the supremum of the naturals. Hence there 
is not a real number \(x\) which is greater than all natural numbers.

\QED

\subsection{Equivalent Formulation}

For any \( \varepsilon > 0 \), there exists \( n \in \N \) such that \( \frac{1}{n} < \varepsilon \)

\subsection{Applications}

\begin{enumerate}
	
	\item \textbf{Density of Rational Numbers:} The Archimedean Principle helps in proving that between 
		   any two real numbers, there exists a rational number.

	\item \textbf{Limits and Infinitesimals:} It ensures that sequences like 
	      \( \left\{ \frac{1}{n} \right\}\) converge to 0, foundational in real analysis and calculus.

	\item \textbf{Bounding Functions:} It is used in analysis to show that functions do not grow faster 
	      than natural numbers in certain contexts.

	\item \textbf{Non-Existence of Infinitely Small Numbers:} The principle implies that real numbers do 
	      not contain infinitesimals (nonzero numbers smaller than all \( \frac{1}{n} \)), distinguishing 
		  \( \R \) from non-standard number systems.

\end{enumerate}

\textbf{Example:}

Prove that \(\forall \varepsilon > 0, \exists N \in \N: \frac{1}{n} < \varepsilon, \forall n \ge N\)

We know that 

\begin{align*}
	\frac{1}{N} &< \varepsilon \\ 
	\varepsilon > \frac{1}{N} &\ge \frac{1}{n} \text{ by transitivity } \\ 
	\varepsilon &> \frac{1}{n}
\end{align*}

We also get \(N > \frac{1}{\varepsilon}\).

And by the Archimedean Principle we let \(b = 1\), \(a = \varepsilon\) and \(n = N\). Then by 
the transitivity we get out proof.

\QED

